\documentclass{article}

% - imports
\usepackage{enumitem}
\usepackage[utf8]{inputenc}
% \usepackage[backend=biber]{biblatex-chicago}
% \addbibresource{bibliography.bib}
\usepackage{url}
\usepackage{parskip}
\usepackage{indentfirst}
% \usepackage{parallel}

\title{\vspace*{-72pt} Concert Report}
\author{Crystal Mandal}
\date{Last Edited \today}

\begin{document}

  \maketitle


For this Concert Report, I attended the ``Tapestry of Sounds'' Concert
as performed by the WashU Symphony Orchestra. I attended the performance
on the evening of the 16th of October.

This concert was presented as a sort of ``overview'' concert, with music
spanning hundreds of years and many styles - in this element, I am a
little disappointed with the results. The concert programme is as
follows:

\begin{enumerate}
\item
  Overture to \emph{Nabucco} by Giuseppe Verdi
\item
  Piano Concerto in A minor Op. 16 by Edvard Grieg
\item
  \emph{A Shelter that Filters the Sun} by Armando Bayolo
\item
  \emph{Soirées Musicales}, Op. 9 by Benjamin Britten.
\end{enumerate}

Though loosely presented as an overview concert - that is, at least in
title and in spirit - much of the musical material centers a 19th
century European soundworld. This sort of unintentional (at least, I hope
it is unintentional) bias/homogeny is something I am quite worried about in
my recitals, and something I'm not sure I'm great at avoiding either. Unfortunately,
it's not really any specific person's fault: Western Art Music glorifies the European
white men of the 18th and 19th centuries so much that any other composers and sounds
get lost.

\end{document}
